% Requires: \usepackage{booktabs} \usepackage{threeparttable}
\begin{table}[htbp]
\centering
\footnotesize
\caption{Performance of each portfolio strategy.}
\label{tab:performance}
\begin{threeparttable}
\begin{tabular}{lrrrrr}
\toprule
 & {Equal-weight} & {EPO} & {PPP} & {Proposed (hist.)} & {Proposed (fwd.)} \\
\midrule
Annualised return (\%) & 11.53 & 7.95 & 10.35 & 29.73 & 35.39 \\
Annualised volatility (\%) & 16.43 & 13.54 & 18.27 & 31.21 & 30.78 \\
Maximum drawdown (\%) & 27.36 & 23.30 & 33.79 & 32.57 & 28.97 \\
Sharpe ratio & 0.63 & 0.49 & 0.53 & 0.93 & 1.08 \\
CRRA certainty equivalent (\%) & 5.33 & 3.95 & 2.03 & 8.07 & 14.04 \\
Annualised turnover & 0.85 & 4.69 & 5.21 & 1.49 & 1.76 \\
Average weight entropy & 5.63 & 4.06 & 4.61 & 0.03 & 0.12 \\
\bottomrule
\end{tabular}
\begin{tablenotes}[flushleft]\footnotesize
\item Sample 2015-02 to 2026-01 (132 monthly observations). Returns are net of 10~bps of one-way transaction cost.
\item The certainty equivalent is the annualised certain return an investor with CRRA utility at $\gamma=5$ would accept in place of the realised series.
\item Turnover is the annualised sum of absolute weight changes. Weight entropy is $-\sum_i w_i \ln w_i$ averaged over rebalance dates; a lower value is a more concentrated book, and $\exp(\cdot)$ of it is an effective number of holdings.
\item Tests of the differences between these strategies and the equal-weight benchmark are reported in Table~\ref{tab:inference}.
\end{tablenotes}
\end{threeparttable}
\end{table}
